\documentclass[11pt]{article}
\usepackage[utf8]{inputenc}
\usepackage{amsmath}
\usepackage{amsfonts}
\usepackage{booktabs}
\usepackage{geometry}
\geometry{a4paper, margin=1in}

\title{The $2^n$ Resonance and Septimal Dark Matter: \\ \large A Treatise on $C=256$ Hz and the 256th-Note Horizon}
\author{Alchemy of Frequency}
\date{March 2026}

\begin{document}

\maketitle

\begin{abstract}
This treatise explores the mathematical convergence of Scientific Pitch ($C=256$ Hz) and hyper-microscopic rhythmic subdivisions. By utilizing 7-limit Just Intonation, we uncover the "Septimal Dark Matter" of the harmonic series—intervals that provide a stable, physical resonance hidden by modern tempering. This document proves that at the 256th-note limit, rhythm and pitch achieve a unified identity.
\end{abstract}

\section{The Physical Prime: $C_4 = 256$ Hz}
We establish Middle C ($C_4$) as exactly $256$ Hz. This aligns the tonal system with binary physics ($2^8$), ensuring that every octave of C is a power of two. This "Scientific Pitch" serves as the $1/1$ ratio, the alchemical source from which all other intervals are distilled \cite{Doty2026}. 

\section{The Alchemical Map: The 5-Limit Lattice}
Just Intonation (JI) is navigated through a tonal lattice rather than a linear scale. Horizontal movement represents the Perfect Fifth ($3:2$), while vertical movement represents the Major Third ($5:4$).

\begin{itemize}
    \item \textbf{Eastward ($3:2$):} From $C$ ($256$ Hz) to $G$ ($384$ Hz).
    \item \textbf{Northward ($5:4$):} From $C$ ($256$ Hz) to $E$ ($320$ Hz).
    \item \textbf{The Synthesis ($15:8$):} The Major Seventh, $B$ ($480$ Hz).
\end{itemize}

\section{Septimal Dark Matter: The 7-Limit}
The "Dark Matter" of the musical universe consists of intervals involving the prime number $7$. These pitches, such as the \textit{Harmonic Seventh} ($7:4$), are physically present in the harmonic series of any string but are discarded by the 12-tone equal temperament (12TET) system \cite{Doty2026}.

\begin{equation}
f_{7} = \frac{7}{4} \cdot 256 = 448 \text{ Hz}
\end{equation}

In 12TET, the minor seventh is tuned to approximately $456.1$ Hz. Our "Dark Matter" $B\flat$ at $448$ Hz is significantly flatter, yet it locks into a "beat-free" stability with Middle C. This creates a sensation of "septimal gravity," a deep resonance that feels grounded rather than restless.

\section{The Horizontal Limit: The 256th Note}
The 256th note (\textit{demisemihemidemisemiquaver}) represents the event horizon of musical duration. In a tempo of 60 BPM (1 Hz), the duration $\Delta t$ is:

\begin{equation}
\Delta t = \frac{1}{64} \text{ s} = 0.015625 \text{ s}
\end{equation}

At $C=256$ Hz, the period of a single vibration is $T = 1/256$ s. Therefore, the 256th note is exactly:
\begin{equation}
\Delta t = 4 \times T
\end{equation}
The rhythmic pulse and the tonal frequency achieve \textit{Phase-Locking}. The distinction between "time" and "pitch" dissolves into a singular auditory event.

\appendix
\section{Frequency Table of the Septimal Scale}
Below is the alchemical distribution of the 7-limit scale, anchored at $C=256$ Hz.

\begin{table}[h!]
\centering
\renewcommand{\arraystretch}{1.5}
\begin{tabular}{@{}llll@{}}
\toprule
Note & Ratio & Frequency (Hz) & Limit Type \\ \midrule
C & $1/1$ & 256.0 & Fundamental \\
D & $9/8$ & 288.0 & 3-Limit \\
E & $5/4$ & 320.0 & 5-Limit \\
F & $4/3$ & 341.33 & 3-Limit \\
G & $3/2$ & 384.0 & 3-Limit \\
A & $5/3$ & 426.66 & 5-Limit \\
B$\flat$ (Septimal) & $7/4$ & 448.0 & 7-Limit (Dark Matter) \\
B & $15/8$ & 480.0 & 5-Limit \\
C' & $2/1$ & 512.0 & Octave \\ \bottomrule
\end{tabular}
\end{table}

\begin{thebibliography}{9}
\bibitem{Doty2026} 
David B. Doty, \textit{The Why, How and Wherefore}, Sound American: The Just Intonation Issue, 2026.
\bibitem{Wiki256} 
Wikipedia Contributors, \textit{Two hundred fifty-sixth note}, Wikipedia, The Free Encyclopedia, 2026.
\end{thebibliography}

\end{document}
